\section{Εισαγωγή}

\subsection{Πλαίσιο}

Στην ενότητα αυτή περιγράφεται το γενικότερο αντικείμενο της εργασίας και
δίνονται οι βασικοί ορισμοί των όρων που χρησιμοποιούνται. Διατυπώνεται το
κύριο πρόβλημα (problem statement) που επιλύεται και αναδεικνύεται η
χρησιμότητα της εργασίας, δηλαδή γιατί το υπό εξέταση ζήτημα είναι
σημαντικό. Ενδεικτική αναφορά σε πηγή της βιβλιογραφίας
\cite{placeholder}. Σύμφωνα με τη βιβλιογραφική ανασκόπηση, το κενό που
καλύπτεται (literature gap) είναι ... .

\subsection{Σκοπός, Ερωτήματα και Στόχοι}

Ο σκοπός (aim) της εργασίας είναι ... . Τα ερευνητικά ερωτήματα (research
questions) που διερευνώνται είναι ... , ενώ οι επιμέρους στόχοι
(objectives) που τίθενται, στο πλαίσιο της εργασίας, είναι ... .

\subsection{Οριοθέτηση}

Στην ενότητα αυτή ορίζεται τι εξαιρείται από τη μελέτη. Περιγράφονται οι
παραδοχές, οι περιορισμοί και το εύρος ισχύος των αποτελεσμάτων ... .

\subsection{Διάρθρωση της Εργασίας}

Η εργασία διαρθρώνεται σε ... κεφάλαια. Στο Κεφάλαιο 1 παρουσιάζεται ... ,
στο Κεφάλαιο 2 ... , και ούτω καθεξής.
